% Lösungen Einheit 5 von 5 – Anwendungen (zusammenbau v0.1)
\einheitenkopf{Einheit 5 von 5 · Anwendungen}

%% TODO Lösung der Erklärzeile 28g) fehlt in der Bank
\erg{28}{a) Anfangswert 45 € ($n = 45$), Änderung 38 € je Stunde ($m = 38$) \quad b) $y = 3x + 8$ \quad c) $y = 0{,}35x + 9$ \quad d) $y = 2x + 4$ \quad e) $y = 40x + 15$ \quad f) $y = 0{,}29x + 1$ \quad g) \ldots \quad h) $m = 6$, $n = 10$: $y = 6x + 10$ \quad i) $85 + 12 \cdot 15 = 265$ € \quad j) A: $20 + 0{,}30 \cdot 60 = 38$ €, B: $0{,}70 \cdot 60 = 42$ €; A ist günstiger \quad k) Angebot 1: $4 \cdot 32 + 120 \cdot 0{,}30 = 164$ €; Angebot 2: $95 + 420 \cdot 0{,}12 = 145{,}40$ €; Angebot 2 ist günstiger; $y = 0{,}12x + 95$}
\erg{29}{a) B (Startwert auf der y-Achse und Anstieg passen)}
\erg{30}{a) zum Beispiel ein Taxi mit 4 € Grundgebühr und 2{,}50 € je km; x km, y Preis in €}
\erg{31}{a) Max hat Anfangswert und Änderung vertauscht; richtig: $m = 2{,}5$, $n = 6$, also $y = 2{,}5x + 6$. \quad b) B: Die Grundgebühr ist gleich, aber jede Minute kostet bei B weniger. \quad c) $n = 12$ (Kosten bei 0 GB), $m = 1{,}5$ (je GB): $y = 1{,}5x + 12$ \quad d) $1{,}9x + 180 = 600$, $x \approx 221$: höchstens 221 km}
